\documentclass[11pt]{article}

\usepackage[a4paper,margin=2.2cm]{geometry}
\usepackage[T1]{fontenc}
\usepackage[utf8]{inputenc}
\usepackage{lmodern}
\usepackage{amsmath,amssymb,amsthm,mathtools}
\usepackage{microtype}
\usepackage{titlesec}
\usepackage[hidelinks]{hyperref}

\titleformat{\section}{\Large\bfseries}{\thesection}{0.75em}{}
\titleformat{\subsection}{\large\bfseries}{\thesubsection}{0.75em}{}
\titlespacing*{\section}{0pt}{1.4ex plus .2ex minus .2ex}{0.9ex}
\titlespacing*{\subsection}{0pt}{1.0ex plus .2ex minus .2ex}{0.5ex}

\newcommand{\R}{\mathbb{R}}

\newtheorem{lemma}{Lemma}[section]
\newtheorem{proposition}[lemma]{Proposition}

% https://claude.ai/chat/f423645c-9097-4641-8a72-944cde37aafc
% https://chatgpt.com/c/69ddf365-c818-8391-bf92-7cd24efe21f7
% https://chatgpt.com/c/69dcd8b5-0f60-8389-8b66-0e61b882405f

\begin{document}
\section{Well-posedness of the LPV-LFR internal loop}

In the LPV-LFR framework of Drenth~\cite{drenth2025thesis}, well-posedness of the
internal interconnection requires that
\begin{equation}
    I - D_{zw}\Delta(Y)
    \quad\text{is nonsingular for all admissible } Y.
    \label{eq:wp_generic}
\end{equation}
We do not invoke Drenth's sufficient condition based on $\rho(D_{zw})<1$.
Instead, we exploit the block structure of the present realization and prove that
\eqref{eq:wp_generic} is equivalent to invertibility of the physical inertia matrix
$M(Y) = M_0 + YM_1 + Y^2M_2$.
Since $Y$ is a physical coordinate with no intrinsic bound, we establish this
equivalence globally for all $Y\in\mathbb{R}$.

\subsection{Preliminary: invertibility of \texorpdfstring{$M_0$}{M0}}

From Proposition~\ref{prop:Minvertible}, $M(Y)\succ 0$ for all $Y\in\mathbb{R}$.
Setting $Y=0$ gives
\begin{equation}
    M_0 = M(0) \succ 0,
    \label{eq:M0pd}
\end{equation}
so $M_0$ is invertible and the matrices $M_0^{-1}$ appearing throughout the
LPV-LFR realization are well-defined.

\subsection{LDU factorization lemma}

The following lemma, which follows directly from block Gaussian elimination
\cite[p.\,1--2]{schurwiki}, is the only tool required for the proof.

\begin{lemma}[LDU factorization]
\label{lem:ldu}
Let
\[
    L =
    \begin{bmatrix}
        A & B \\ C & D
    \end{bmatrix}
\]
be a block matrix with $D$ invertible.  Define the Schur complement of $D$ in
$L$ as
\begin{equation}
    S := A - BD^{-1}C.
    \label{eq:schur_def}
\end{equation}
Then $L$ admits the factorization
\begin{equation}
    L
    =
    \underbrace{
    \begin{bmatrix}
        I & BD^{-1} \\ 0 & I
    \end{bmatrix}
    }_{U}
    \underbrace{
    \begin{bmatrix}
        S & 0 \\ 0 & D
    \end{bmatrix}
    }_{\Lambda}
    \underbrace{
    \begin{bmatrix}
        I & 0 \\ D^{-1}C & I
    \end{bmatrix}
    }_{V}.
    \label{eq:ldu}
\end{equation}
Since $U$ and $V$ are unit triangular, $\det U = \det V = 1$, and therefore
\begin{equation}
    \det L = \det(D)\,\det(S).
    \label{eq:detschur}
\end{equation}
Consequently, $L$ is invertible if and only if both $D$ and $S$ are invertible.
\end{lemma}

\begin{proof}
Direct multiplication of the right-hand side of \eqref{eq:ldu}:
\[
    U\Lambda V
    =
    \begin{bmatrix}
        S + BD^{-1}C & D \cdot D^{-1}C \cdot 0 + \ldots
    \end{bmatrix}.
\]
Computing each block explicitly:
\begin{align*}
    (U\Lambda V)_{11}
        &= S + BD^{-1}C = (A - BD^{-1}C) + BD^{-1}C = A,\\
    (U\Lambda V)_{12}
        &= BD^{-1}D = B,\\
    (U\Lambda V)_{21}
        &= D\cdot D^{-1}C = C,\\
    (U\Lambda V)_{22}
        &= D,
\end{align*}
which recovers $L$ exactly.
The determinant relation \eqref{eq:detschur} follows because
$\det L = \det(U)\det(\Lambda)\det(V) = 1\cdot\det(D)\det(S)\cdot 1$.
\end{proof}

\subsection{Proposition: plant-specific equivalence}

\begin{proposition}
\label{prop:wellposed}
For the gantry LPV-LFR realization with
\begin{equation}
    D_{zw} =
    \begin{bmatrix}
        -M_0^{-1}M_1 & -M_0^{-1}M_2 \\
        I_3 & 0
    \end{bmatrix},
    \qquad
    \Delta(Y) = YI_6,
    \label{eq:DzwDelta}
\end{equation}
the following equivalence holds for all $Y\in\mathbb{R}$:
\begin{equation}
    I - D_{zw}\Delta(Y)\ \text{is invertible}
    \quad\Longleftrightarrow\quad
    M(Y)\ \text{is invertible}.
    \label{eq:mainequiv}
\end{equation}
\end{proposition}

\begin{proof}
Define $L(Y) := I_6 - D_{zw}\Delta(Y)$.
Since $\Delta(Y) = YI_6$, direct multiplication gives
\begin{equation}
    D_{zw}\Delta(Y)
    =
    \begin{bmatrix}
        -YM_0^{-1}M_1 & -YM_0^{-1}M_2 \\
        YI_3 & 0
    \end{bmatrix},
\end{equation}
hence
\begin{equation}
    L(Y)
    =
    \begin{bmatrix}
        I_3 + YM_0^{-1}M_1 & YM_0^{-1}M_2 \\
        -YI_3 & I_3
    \end{bmatrix}
    =:
    \begin{bmatrix}
        A(Y) & B(Y) \\ C(Y) & D(Y)
    \end{bmatrix}.
    \label{eq:LY}
\end{equation}
The lower-right block is $D(Y) = I_3$, which is invertible for all $Y\in\mathbb{R}$
and is therefore the natural pivot for Lemma~\ref{lem:ldu}.
The upper-left block $A(Y) = I_3 + YM_0^{-1}M_1$ has $Y$-dependent invertibility
and cannot serve as pivot without additional restriction on $Y$.

Since $D(Y) = I_3$ is invertible, Lemma~\ref{lem:ldu} applies and gives:
$L(Y)$ is invertible if and only if the Schur complement
\begin{equation}
    S(Y) := A(Y) - B(Y)D(Y)^{-1}C(Y)
\end{equation}
is invertible.  Computing $S(Y)$ with $D(Y)^{-1} = I_3$:
\begin{align}
    S(Y)
    &= \bigl(I_3 + YM_0^{-1}M_1\bigr)
       - \bigl(YM_0^{-1}M_2\bigr)\bigl(-YI_3\bigr) \notag\\
    &= I_3 + YM_0^{-1}M_1 + Y^2M_0^{-1}M_2 \notag\\
    &= M_0^{-1}\bigl(M_0 + YM_1 + Y^2M_2\bigr)
       \qquad\text{(factoring } M_0^{-1}\text{ from each term)} \notag\\
    &= M_0^{-1}M(Y).
    \label{eq:SY}
\end{align}
Since $M_0$ is invertible by \eqref{eq:M0pd}, $\det(M_0^{-1})\neq 0$.
From the determinant relation \eqref{eq:detschur},
\begin{equation}
    \det S(Y) = \det(M_0^{-1})\,\det(M(Y)),
\end{equation}
so $\det S(Y) = 0$ if and only if $\det M(Y) = 0$.  Therefore
\begin{equation}
    S(Y)\ \text{is invertible}
    \quad\Longleftrightarrow\quad
    M(Y)\ \text{is invertible}.
\end{equation}
Combining with Lemma~\ref{lem:ldu} proves \eqref{eq:mainequiv}.
\end{proof}

\subsection{Global well-posedness}

The internal loop equation of the LPV-LFR is
\begin{equation}
    \bigl(I - D_{zw}\Delta(Y)\bigr)\,z = C_z x + D_{zu} u.
    \label{eq:loopequation}
\end{equation}
By Proposition~\ref{prop:wellposed}, this has a unique solution in $z$ for every
$x$, $u$, and $Y\in\mathbb{R}$ if and only if $M(Y)$ is invertible for all
$Y\in\mathbb{R}$.  Since $M(Y)\succ 0$ for all $Y\in\mathbb{R}$ by
Proposition~\ref{prop:Minvertible}, the realization is globally well-posed.

This is a plant-specific proof of Drenth's necessary-and-sufficient
well-posedness condition~\eqref{eq:wp_generic}.  Drenth's sufficient
spectral-radius theorem is not invoked.
% PASTE SECTION HERE

\newpage

\begin{thebibliography}{9}

\bibitem{drenth2025thesis}
R.~Drenth,
``Efficient Gradient-Based Learning of {LPV} Models with Linear Fractional Representation,''
M.Sc.\ thesis, Eindhoven University of Technology, Sept.\ 2025.

\bibitem{schurwiki}
Wikipedia contributors,
``Schur complement,''
\textit{Wikipedia, The Free Encyclopedia},
2024.
\url{https://en.wikipedia.org/wiki/Schur_complement}

\end{thebibliography}

\end{document}

The best match for your derivation is the dedicated monograph cited in the Wikipedia PDF you already uploaded:
Zhang, F. (Ed.) (2005). The Schur Complement and Its Applications. Numerical Methods and Algorithms, Vol. 4. Springer.
This is the standard academic reference for exactly the LDU factorization, determinant formula, and invertibility equivalence your proof uses. The Wikipedia article you uploaded cites it directly as reference [2], so the chain is: your proof → Wikipedia PDF → Zhang (2005). Using Zhang (2005) directly is cleaner than citing a Wikipedia article in a thesis.
The second viable option is:
Horn, R.A. and Johnson, C.R. (2013). Matrix Analysis, 2nd ed. Cambridge University Press. §0.8.
Horn & Johnson covers block matrix LDU factorization and the Schur complement as part of standard matrix theory, and you already cite it in Document 2 for Sylvester's criterion. Using the same source in both companion notes would be internally consistent.
For completeness, Boyd & Vandenberghe, Convex Optimization (2004), Appendix A.5.5 also states the result, but only for symmetric matrices, so it does not cover your general non-symmetric case and should not be used here.
Recommendation: cite Zhang (2005) for the Schur complement lemma in your well-posedness section, and replace \cite{schurwiki} in the wrapper with:
@book{zhang2005schur,
  author    = {Zhang, Fuzhen},
  title     = {The {Schur} Complement and Its Applications},
  series    = {Numerical Methods and Algorithms},
  volume    = {4},
  publisher = {Springer},
  year      = {2005},
  doi       = {10.1007/b105056}
}